\documentclass[11pt,a4paper]{article}
\usepackage[T1]{fontenc}
\usepackage{lmodern}
\usepackage[margin=27mm]{geometry}
\usepackage{amsmath,amssymb,amsthm,mathtools}
\usepackage{microtype}
\usepackage{booktabs,array}
\usepackage{enumitem}
\usepackage[hidelinks]{hyperref}
\setlength{\parskip}{5pt plus 1pt}
\setlength{\parindent}{0pt}
\setlength{\emergencystretch}{3em}
\setlist{itemsep=2pt,topsep=3pt}
\raggedbottom
\newtheorem{theorem}{Theorem}[section]
\newtheorem{lemma}[theorem]{Lemma}
\newtheorem{corollary}[theorem]{Corollary}
\newtheorem{proposition}[theorem]{Proposition}
\theoremstyle{definition}
\newtheorem{definition}[theorem]{Definition}
\newtheorem{example}[theorem]{Example}
\theoremstyle{remark}
\newtheorem{remark}[theorem]{Remark}
\newcommand{\Z}{\mathbb{Z}}
\newcommand{\N}{\mathbb{N}}
\newcommand{\Q}{\mathbb{Q}}
\newcommand{\F}{\mathbb{F}}
\newcommand{\ps}[1]{#1[[x]]}
\newcommand{\coeff}[2]{[x^{#1}]\,#2}
\newcommand{\hl}[1]{\textbf{#1}}
\newcommand{\para}[1]{\par\medskip\noindent\textbf{#1}\enspace}
\hypersetup{pdftitle={A greedy sequence that avoids sums of three consecutive terms},pdfauthor={Dev D. Goyal},pdfsubject={OEIS A394576}}
\title{A greedy sequence that avoids sums of three consecutive terms\\[6pt]\large A proof of the first conjecture in OEIS A394576}
\author{Dev D. Goyal}
\date{October 2026}
\begin{document}
\maketitle
\begin{abstract}
The sequence A394576 of the On-Line Encyclopedia of Integer Sequences (OEIS), contributed by Hoang Xuan Thanh, starts
with $a(1)=1$ and then always takes the smallest integer larger than the previous term that is not a sum
$a(k)+a(k+1)+a(k+2)$ of three consecutive earlier terms. Its first conjecture is the bound
$a(n)\le\lfloor(4n-1)/3\rfloor$. We prove this bound, together with the lower bound
$a(n)\ge\lfloor(4n-3)/3\rfloor$, for every $n\ge1$. The two bounds differ by at most one and coincide when $3\mid n$, so
$a(n)$ is always the upper bound or one less than it, and $a(3k)=4k-1$. The proof rests on an explicit construction: for
each $k\ge1$, exactly one of the four integers $4k,4k+1,4k+2,4k+3$ is a sum of three consecutive terms, and the other
three are the terms $a(3k+1)$, $a(3k+2)$, $a(3k+3)$. We prove that the sequence so constructed obeys the greedy rule;
since the rule determines the sequence uniquely, it is the sequence of the entry. The results have been checked by
computer and proved formally in the Lean proof assistant.
\end{abstract}
\begin{center}\begin{minipage}{0.88\linewidth}\small
\textbf{How this paper was produced.} This paper reports the result of an AI research effort directed by
Dev~D.~Goyal. ChatGPT (OpenAI), using the GPT-6 model with Pro reasoning and working autonomously from prompts by
Dev~D.~Goyal, found the proof, wrote its Lean~4 formalization, and wrote a first draft of the paper.
Dev~D.~Goyal chose the problems and asked for proofs, formal verification and revisions, but did not supply proof
ideas. Claude Opus~5.5, used with Claude Code, then checked the work independently: it rebuilt the formal proofs from
source and confirmed that they rely only on Lean's three standard axioms, recomputed the sequence from its definition
compared it with the published terms and tested every claim numerically, and read every step of the proofs. It also wrote the present
version of this paper.
Dev~D.~Goyal has reviewed the proofs and is responsible for them; this review is not an independent expert validation.
The paper has not been peer reviewed.
The section ``About this paper'' at the end gives the details.
\end{minipage}\end{center}

\section{Introduction}

\subsection{The sequence}\label{sec:sequence}
The \emph{On-Line Encyclopedia of Integer Sequences} (OEIS)~\cite{oeisdb} is a public database of integer sequences.
Many entries record \emph{conjectures}: patterns observed in the data that the contributor did not prove. This paper
concerns entry A394576~\cite{oeis}, contributed by Hoang Xuan Thanh in 2026, which reads in part:
\begin{quote}\small\ttfamily\raggedright a(1) = 1; a(n+1) is the smallest positive integer greater than a(n) that avoids violating the condition a(k) + a(k+1) + a(k+2) is not a term.\par\smallskip Conjecture 1: a(n) <= floor((4*n-1)/3).\par\smallskip\hfill\textrm{--- OEIS A394576, Hoang Xuan Thanh, May 2026}\end{quote}

In other words, the sequence is built one term at a time, and the rule is that no sum of three consecutive terms may itself be
a term. When $a(n+1)$ is chosen, the sums $a(k)+a(k+1)+a(k+2)$ with $k+2\le n$ are already known, and $a(n+1)$ must avoid
them. A sum that involves $a(n+1)$ is larger than $a(n+1)$, so it constrains only later terms.
Precisely, the rule is: $a(1)=1$, and for $n\ge1$,
\begin{equation}\label{eq:rule}
a(n+1)=\min\bigl\{t>a(n):\ t\ne a(k)+a(k+1)+a(k+2)\ \text{for all }k\ge1\text{ with }k+2\le n\bigr\}.
\end{equation}
The first terms are 1, 2, 3, 4, 5, 7, 8, 10, 11, 13, 14, 15, 17, 18, 19, 21, 22, 23, 24, 26, \dots{} For instance, $a(6)=7$ because
$a(1)+a(2)+a(3)=6$ is forbidden.

\subsection{The question}
The first conjecture of the entry says that $a(n)\le\lfloor(4n-1)/3\rfloor$, where $\lfloor y\rfloor$ denotes the
largest integer not exceeding $y$. The table compares the first $32$ terms with this bound and with the lower bound
$\lfloor(4n-3)/3\rfloor$ proved below. The indices in bold are the multiples of $3$, where the two bounds coincide.
\begin{center}\small\setlength{\tabcolsep}{3.2pt}
\begin{tabular}{>{\raggedright\arraybackslash}p{4.6em}>{\centering\arraybackslash}p{1.15em}>{\centering\arraybackslash}p{1.15em}>{\centering\arraybackslash}p{1.15em}>{\centering\arraybackslash}p{1.15em}>{\centering\arraybackslash}p{1.15em}>{\centering\arraybackslash}p{1.15em}>{\centering\arraybackslash}p{1.15em}>{\centering\arraybackslash}p{1.15em}>{\centering\arraybackslash}p{1.15em}>{\centering\arraybackslash}p{1.15em}>{\centering\arraybackslash}p{1.15em}>{\centering\arraybackslash}p{1.15em}>{\centering\arraybackslash}p{1.15em}>{\centering\arraybackslash}p{1.15em}>{\centering\arraybackslash}p{1.15em}>{\centering\arraybackslash}p{1.15em}}
\toprule
$n$ & 1 & 2 & \hl{3} & 4 & 5 & \hl{6} & 7 & 8 & \hl{9} & 10 & 11 & \hl{12} & 13 & 14 & \hl{15} & 16 \\ \midrule
$\lfloor\frac{4n-3}3\rfloor$ & 0 & 1 & \hl{3} & 4 & 5 & \hl{7} & 8 & 9 & \hl{11} & 12 & 13 & \hl{15} & 16 & 17 & \hl{19} & 20 \\
$a(n)$ & 1 & 2 & \hl{3} & 4 & 5 & \hl{7} & 8 & 10 & \hl{11} & 13 & 14 & \hl{15} & 17 & 18 & \hl{19} & 21 \\
$\lfloor\frac{4n-1}3\rfloor$ & 1 & 2 & \hl{3} & 5 & 6 & \hl{7} & 9 & 10 & \hl{11} & 13 & 14 & \hl{15} & 17 & 18 & \hl{19} & 21 \\
\bottomrule
\end{tabular}
\\[4pt]
\begin{tabular}{>{\raggedright\arraybackslash}p{4.6em}>{\centering\arraybackslash}p{1.15em}>{\centering\arraybackslash}p{1.15em}>{\centering\arraybackslash}p{1.15em}>{\centering\arraybackslash}p{1.15em}>{\centering\arraybackslash}p{1.15em}>{\centering\arraybackslash}p{1.15em}>{\centering\arraybackslash}p{1.15em}>{\centering\arraybackslash}p{1.15em}>{\centering\arraybackslash}p{1.15em}>{\centering\arraybackslash}p{1.15em}>{\centering\arraybackslash}p{1.15em}>{\centering\arraybackslash}p{1.15em}>{\centering\arraybackslash}p{1.15em}>{\centering\arraybackslash}p{1.15em}>{\centering\arraybackslash}p{1.15em}>{\centering\arraybackslash}p{1.15em}}
\toprule
$n$ & 17 & \hl{18} & 19 & 20 & \hl{21} & 22 & 23 & \hl{24} & 25 & 26 & \hl{27} & 28 & 29 & \hl{30} & 31 & 32 \\ \midrule
$\lfloor\frac{4n-3}3\rfloor$ & 21 & \hl{23} & 24 & 25 & \hl{27} & 28 & 29 & \hl{31} & 32 & 33 & \hl{35} & 36 & 37 & \hl{39} & 40 & 41 \\
$a(n)$ & 22 & \hl{23} & 24 & 26 & \hl{27} & 28 & 30 & \hl{31} & 32 & 33 & \hl{35} & 36 & 37 & \hl{39} & 40 & 41 \\
$\lfloor\frac{4n-1}3\rfloor$ & 22 & \hl{23} & 25 & 26 & \hl{27} & 29 & 30 & \hl{31} & 33 & 34 & \hl{35} & 37 & 38 & \hl{39} & 41 & 42 \\
\bottomrule
\end{tabular}
\end{center}

The upper bound $\lfloor(4n-1)/3\rfloor$ is the $n$-th positive integer not divisible by $4$ (OEIS A042968): the
sequence follows the numbers $1,2,3,5,6,7,9,\dots$ closely, but occasionally falls one step behind.

\subsection{What this paper proves}
\textbf{Main result} (Theorem~\ref{thm:main}). \emph{For every $n\ge1$,}
\[
\Bigl\lfloor\frac{4n-3}{3}\Bigr\rfloor\le a(n)\le\Bigl\lfloor\frac{4n-1}{3}\Bigr\rfloor .
\]
The upper bound is Conjecture~1 of the entry. The two bounds differ by at most $1$ and agree when $n$ is a multiple of
$3$, so $a(3k)=4k-1$ for every $k\ge1$ (Corollary~\ref{cor:error}).

\subsection{What was known}
The entry records two further conjectures about where $a(n)$ equals the upper bound and where it is one less.
Thanh revised Conjecture~2 and added Conjecture~3 in September 2026. Conjecture~2 states that the indices $n$ with $a(n)=\lfloor(4n-1)/3\rfloor$
include runs, which the entry calls stable blocks, of lengths $9^{k-1}+2$ at explicitly given positions, separated by
gaps. Conjecture~3 describes the gaps: writing $d(n)=\lfloor(4n-1)/3\rfloor-a(n)$, $X=(1,0,0)$ and $Y=(1,1,0)$, it states
that the values of $d(n)$ in the $k$-th gap form the word $X^{A_k}Y^{B_k}X^{C_k}(1)$, for explicit exponents $A_k$,
$B_k$, $C_k$. These two conjectures are not proved here. Thanh has deposited software that computes the sequence and its
generalizations~\cite{algorithm}. When the entry was checked on 24~October~2026, the bound was still labelled as a
conjecture, and we are not aware of an earlier proof.

\subsection{Motivation}
Greedy sequences that avoid a forbidden pattern are a source of difficult problems: the rule is simple, but the
global behavior of the sequence is often irregular and difficult to predict. Here the behavior is very
regular: for each $k\ge1$, exactly one of the four integers $4k,\dots,4k+3$ is a forbidden sum and the other three are
terms, so the sequence has density $3/4$ in the positive integers. The proof shows how a greedy definition can
be replaced by an explicit construction, and how to prove that the two agree.

\subsection{How to read this paper}
Section~\ref{sec:background} fixes the notation. Section~\ref{sec:results} states the results.
Section~\ref{sec:proofs} gives the construction and proves that it yields the greedy sequence.
Section~\ref{sec:checks} describes the computer checks and the formal proof, and Section~\ref{sec:discussion} discusses
the remaining conjectures.
\section{Background}\label{sec:background}
\subsection{Notation}
For a real number $y$, $\lfloor y\rfloor$ is the largest integer not exceeding $y$. We write
\[
L(n)=\Bigl\lfloor\frac{4n-3}{3}\Bigr\rfloor,\qquad U(n)=\Bigl\lfloor\frac{4n-1}{3}\Bigr\rfloor\qquad(n\ge1).
\]
Writing $n=3k+r$ with $r\in\{1,2,3\}$, one finds
\begin{equation}\label{eq:LU}
\begin{array}{c|ccc}
n&3k+1&3k+2&3k+3\\ \hline
L(n)&4k&4k+1&4k+3\\
U(n)&4k+1&4k+2&4k+3
\end{array}
\end{equation}
for every $k\ge0$. In particular $0\le U(n)-L(n)\le1$, and $L(n)=U(n)$ when $3\mid n$. For $k\ge1$ we
write $I_k=\{4k,4k+1,4k+2,4k+3\}$; these intervals are disjoint, and together they cover the integers from $4$ on.

\subsection{Greedy definitions}
A sequence defined by a rule of the form \eqref{eq:rule}, in which each new term is the least admissible integer, is
called \emph{greedy}. At each stage only finitely many sums are forbidden, so the minimum in \eqref{eq:rule} exists,
and the rule defines a sequence. It defines only one: two sequences that satisfy \eqref{eq:rule} and agree up to index
$n$ have the same forbidden sums at that stage, hence the same next term, and by induction they agree everywhere. The
difficulty in proving properties of greedy sequences is that the rule is local, while statements such as a bound for
$a(n)$ are global. The proof below replaces the rule by an explicit construction.

\section{Results}\label{sec:results}

\begin{theorem}\label{thm:main}
The rule \eqref{eq:rule} has exactly one solution, and for every $n\ge1$,
\[
L(n)\le a(n)\le U(n).
\]
Moreover, the terms of the sequence and the sums $a(k)+a(k+1)+a(k+2)$, $k\ge1$, are disjoint sets whose union is the set
of all positive integers.
\end{theorem}

\begin{corollary}\label{cor:error}
For every $n\ge1$, $U(n)-a(n)\in\{0,1\}$; in particular Conjecture~1 of the entry holds. For every $k\ge1$,
$a(3k)=4k-1$.
\end{corollary}

\begin{proof}
This follows from Theorem~\ref{thm:main} and the table \eqref{eq:LU}.
\end{proof}

\section{Proofs}\label{sec:proofs}

\para{The idea in brief.} We construct a sequence $A$ explicitly, three terms at a time, and then show that $A=a$. The
first three terms are $1,2,3$. At stage $k\ge1$ the sum $b(k)=A(k)+A(k+1)+A(k+2)$ is already known, and the next three
terms are the elements of $I_k$ other than $b(k)$. The bounds $L\le A\le U$ show that $b(k)$ always lies in $I_k$, so
this is always possible. Finally, every integer that the construction skips is one of the sums $b(k)$, which shows that
$A$ obeys the greedy rule.

\subsection{The construction}
Put $A(1)=1$, $A(2)=2$, $A(3)=3$. Suppose $A(1),\dots,A(3k)$ have been defined, with $k\ge1$. Since $k+2\le3k$, the number
\begin{equation}\label{eq:b}
b(k)=A(k)+A(k+1)+A(k+2)
\end{equation}
is known. Let $A(3k+1)<A(3k+2)<A(3k+3)$ be the elements of $I_k\setminus\{b(k)\}$, listed in increasing order (if $b(k)\notin
I_k$, take the three smallest elements of $I_k$; Lemma~\ref{lem:bounds} shows that this case never occurs). For example,
$b(1)=6$ gives the terms $4,5,7$, and then $b(2)=2+3+4=9$ gives $8,10,11$.

\begin{lemma}\label{lem:sums}
For every $k\ge1$, $L(k)+L(k+1)+L(k+2)=4k$ and $U(k)+U(k+1)+U(k+2)=4k+2$.
\end{lemma}

\begin{proof}
The three numbers $4k-3$, $4k+1$, $4k+5$ have the three different remainders $0,1,2$ modulo $3$, in some order.
Subtracting these remainders and dividing by $3$ gives $L(k)+L(k+1)+L(k+2)=\frac{(12k+3)-3}{3}=4k$. The same argument
with $4k-1$, $4k+3$, $4k+7$ gives $U(k)+U(k+1)+U(k+2)=\frac{(12k+9)-3}{3}=4k+2$.
\end{proof}

\begin{lemma}\label{lem:bounds}
For every $n\ge1$, $L(n)\le A(n)\le U(n)$, and for every $k\ge1$, $4k\le b(k)\le4k+2$.
\end{lemma}

\begin{proof}
The bounds hold for $n=1,2,3$ by \eqref{eq:LU}. Suppose they hold up to $n=3k$. Since $k+2\le3k$, Lemma~\ref{lem:sums}
gives $4k\le b(k)\le4k+2$. Write $b(k)=4k+s$ with $s\in\{0,1,2\}$. Removing $4k+s$ from $\{4k,\dots,4k+3\}$ leaves a
smallest element in $\{4k,4k+1\}$, a middle element in $\{4k+1,4k+2\}$, and the largest element $4k+3$. By
\eqref{eq:LU} these sets are exactly the ranges $[L(n),U(n)]$ for $n=3k+1,3k+2,3k+3$. This completes the induction.
\end{proof}

\begin{lemma}\label{lem:partition}
The sequence $A$ is strictly increasing. Its terms and the numbers $b(k)$, $k\ge1$, form two disjoint sets whose union is
the set of positive integers.
\end{lemma}

\begin{proof}
The first three terms are $1,2,3$. By Lemma~\ref{lem:bounds}, $b(k)\in I_k$, and the terms of stage $k$ are the other
three elements of $I_k$, in increasing order. Since the intervals $I_k$ are disjoint and cover the integers from $4$ on,
the claims follow.
\end{proof}

\begin{proof}[Proof of Theorem~\ref{thm:main}]
We show that $A$ satisfies \eqref{eq:rule}. Fix $n\ge1$. First, $A(n+1)>A(n)$, and $A(n+1)$ is not equal to any $b(k)$,
by Lemma~\ref{lem:partition}; so $A(n+1)$ is admissible. Next, let $t$ be an integer with $A(n)<t<A(n+1)$. By
Lemma~\ref{lem:partition}, $t=b(k)$ for some $k\ge1$. Since all terms are positive, $A(k+2)<b(k)=t<A(n+1)$, and since $A$
is increasing, $k+2<n+1$, that is $k+2\le n$. So $t$ is one of the sums forbidden at stage $n$, and $A(n+1)$ is the least
admissible integer. Hence $A$ satisfies \eqref{eq:rule}. As explained in Section~\ref{sec:background}, \eqref{eq:rule}
has only one solution, so $A=a$. The bounds now follow from Lemma~\ref{lem:bounds}, and the last statement from
Lemma~\ref{lem:partition}.
\end{proof}

\section{Checking the results}\label{sec:checks}

A mathematical proof can contain mistakes, so the results were checked in two independent ways.

\para{Numerical check.}
We computed the sequence directly from the greedy rule \eqref{eq:rule}, for $n\le100{,}000$. This reproduces all
75 terms printed in the OEIS entry, and confirms the bounds of Theorem~\ref{thm:main}, the values $a(3k)=4k-1$, and that
no term is a sum of three consecutive terms, throughout that range.
The computation was done with the PARI/GP computer algebra system, using code that was written separately from the
formal proofs (file \texttt{verification/independent\_checks.gp} in the accompanying repository~\cite{repo}).
A computation can only test finitely many cases, so it is a check on the definitions and on the indexing, not a proof.

\para{Formal proof in Lean.}
Lean~\cite{lean} is a \emph{proof assistant}: a program in which mathematical statements and their proofs are written
in a precise formal language, and every step is checked mechanically down to the basic rules of logic. The theorems of Section~\ref{sec:results} are proved in Lean~4 (version 4.35.0) on top of
Mathlib~\cite{mathlib}, Lean's library of formalized mathematics. The statements are collected in the file \texttt{lean/Entries/A394576.lean} of the repository,
in the namespace \texttt{OEIS.A394576}. There the sequence is indexed from $0$: \texttt{sequence n} is $a(n+1)$, and the rule is written out
explicitly in the statement of \texttt{source\_existence\_uniqueness}. The construction and its properties are in
\texttt{lean/Families/GreedyTriple.lean} (namespace \texttt{Families.GreedyTriple}).
The table lists the results covered by the Lean formalization. The Lean formalization does not cover the additional arguments and observations in Section~\ref{sec:discussion}.
\begin{center}\small\renewcommand{\arraystretch}{1.25}
\begin{tabular}{>{\raggedright\arraybackslash}p{0.52\linewidth}>{\raggedright\arraybackslash}p{0.4\linewidth}}
\toprule
Statement in this paper & Lean theorem(s)\\
\midrule
The rule \eqref{eq:rule} has exactly one solution & \texttt{source\_\allowbreak{}existence\_\allowbreak{}uniqueness} \\
The construction satisfies \eqref{eq:rule} & \texttt{canonical\_\allowbreak{}source} \\
Theorem~\ref{thm:main}, the bounds & \texttt{canonical\_\allowbreak{}bounds}, \texttt{source\_\allowbreak{}implies\_\allowbreak{}bound} \\
Theorem~\ref{thm:main}, the partition of the positive integers & \texttt{complementary\_\allowbreak{}ranges} \\
Corollary~\ref{cor:error} & \texttt{bounded\_\allowbreak{}error}, \texttt{exact\_\allowbreak{}every\_\allowbreak{}third} \\
\bottomrule
\end{tabular}
\end{center}
The Lean project was rebuilt from source, and the axioms on which each of these theorems depends were listed with
Lean's \texttt{\#print axioms} command. Every one depends only on the three standard axioms of Lean's logic
(\texttt{propext}, \texttt{Classical.choice}, \texttt{Quot.sound}); no step is left unproved. A formal proof certifies
exactly the statement that Lean checks. Whether that statement says what the OEIS entry means is a separate question,
so the Lean definitions were also compared with the definition in the OEIS entry
(Section~\ref{sec:sequence}). The principal theorems are stated once more, using only Mathlib's definitions, in the
file \texttt{lean/Challenges/FurtherResults.lean}, and Lean's Comparator tool confirms that the
proved theorems have exactly these statements.

\section{Discussion}\label{sec:discussion}

\para{Scope.} Theorem~\ref{thm:main} proves Conjecture~1 of the entry and a matching lower bound. It does not prove
Conjectures~2 and~3, which describe exactly where $a(n)=U(n)$ and where $a(n)=U(n)-1$. In the language of this paper,
they describe the position $s=b(k)-4k\in\{0,1,2\}$ of the forbidden sum in each interval $I_k$: by the proof of
Lemma~\ref{lem:bounds}, the differences $U(n)-a(n)$ for $n=3k+1,3k+2,3k+3$ are $(0,0,0)$, $(1,0,0)$ or $(1,1,0)$ according as
$s=0$, $1$ or $2$. These are the word $(0,0,0)$ and the words $X$ and $Y$ of Conjecture~3, so the conjectures amount to
a description of the sequence $s(k)=b(k)-4k$.

\para{The skipped numbers.} By Theorem~\ref{thm:main}, the positive integers not in the sequence are exactly the sums
$b(k)=a(k)+a(k+1)+a(k+2)$, $k\ge1$: $6,9,12,16,20,\dots$, one in each interval $I_k$.

\para{Sums of more terms.} The entry refers to the analogous sequence A394676, in which sums of four consecutive terms
are forbidden. An argument of the same kind would need intervals that each contain exactly one forbidden sum; it would
be interesting to know whether such a structure exists there.

\section*{About this paper}\label{sec:about}
\para{How the work was done.}
This paper is one of 28 papers on conjectures from the OEIS that came out of an AI research effort directed by
Dev~D.~Goyal~\cite{repo}. The roles were as follows.
\begin{itemize}
\item \emph{Dev~D.~Goyal} started the project, chose which problems to pursue, asked for proofs, formal verification and
revisions, and decided what to publish. Dev~D.~Goyal did not supply proof strategies or intermediate arguments.
\item \emph{ChatGPT} (OpenAI), using the GPT-6 model with Pro reasoning, worked autonomously from those prompts. It
searched the OEIS for suitable conjectures, found the proofs, wrote the Lean~4 formalization, and wrote a first draft
of this paper.
\item \emph{Claude Opus~5.5}, used with Claude Code, checked the work independently. It rebuilt the Lean project from source and confirmed that the theorems listed in Section~\ref{sec:checks} rely only on
Lean's standard axioms; compared the Lean definitions with the OEIS definition; recomputed the sequence with separately
written PARI/GP code, compared it with the terms published in the OEIS, and tested every claim of the paper on many
terms; and read every step of the proofs. It then wrote the present version of the paper, expanding the
explanations and, in places, writing out an argument in more detail or by a more elementary route. The theorems of
Section~\ref{sec:results} and the proof strategies are ChatGPT's. The Discussion section, including the short arguments
in it, was written by Claude Opus~5.5; its observations were checked by computer and are not part of the Lean
formalization.
\end{itemize}
Dev~D.~Goyal has reviewed the proofs and is responsible for this paper; this review is not an independent expert
validation.

\para{Review status.} This paper has not been peer reviewed, and it has not been submitted to a journal or to the OEIS.

\para{Priority.} An OEIS entry that labels a statement as a conjecture shows only that its contributor did not know
a proof; it does not show that no proof existed. During the research only limited searches for earlier work were
made. No claim of priority is made.

\para{Corrections.} Please report any error in the issue tracker of the repository~\cite{repo}, where the Lean sources and the numerical checks are also
available.

\section*{Acknowledgements}
The sequence and its conjectures are due to Hoang~Xuan~Thanh~\cite{oeis}, who also deposited software that computes the
sequence~\cite{algorithm}.

\begin{thebibliography}{99}\small
\raggedright
\bibitem{oeis} H.~X.~Thanh, entry A394576 ($a(n+1)$ is the smallest integer greater than $a(n)$ avoiding sums of three consecutive terms), \emph{The On-Line Encyclopedia of Integer Sequences}, May 2026. \url{https://oeis.org/A394576}. Text as of October 2026.
\bibitem{algorithm} H.~X.~Thanh, \emph{Algorithm for OEIS A394576, A394676 and generalized forbidden sums}, software deposit, Zenodo, 2026. \url{https://doi.org/10.5281/zenodo.22674326}.
\bibitem{oeisdb} N.~J.~A.~Sloane and The OEIS Foundation Inc., \emph{The On-Line Encyclopedia of Integer Sequences}. \url{https://oeis.org}.
\bibitem{lean} L.~de~Moura and S.~Ullrich, The Lean~4 theorem prover and programming language, in \emph{Automated Deduction -- CADE 28}, Lecture Notes in Computer Science 12699, Springer, 2021, 625--635.
\bibitem{mathlib} The mathlib Community, The Lean mathematical library, in \emph{Proceedings of the 9th ACM SIGPLAN International Conference on Certified Programs and Proofs (CPP 2020)}, 367--381. Version used: commit \texttt{56a005855ea1}.
\bibitem{repo} D.~D.~Goyal, \emph{Proofs and corrections of OEIS conjectures}, repository, 2026. \url{https://github.com/dvgyl/oeis-proofs}.
\end{thebibliography}
\end{document}
